\documentclass[UTF8]{ctexart}

% 支持从项目根目录、学科目录或当前文档目录读取共享样式。
\makeatletter
\def\input@path{{../}{../../}}
\makeatother
\usepackage{researchnotes}

\title{概率论与统计学的区别与联系}
\author{}
\date{}

\begin{document}
\maketitle

\section{研究同类问题，但出发点不同}
\label{sec:prob-stat-distinction}

概率论与统计学都关心随机现象和不确定性，但二者不是同一个概念。
\keyterm{概率论}（probability theory）主要研究给定数学模型下随机现象的规律，
例如事件发生的概率、随机变量的分布，以及大量重复试验中呈现的规律。
\keyterm{统计学}（statistics）主要研究如何收集和利用数据，
描述现象、推断未知规律并作出决策，同时分析结论的不确定性。
本文通过抛硬币这一共同问题说明二者的区别，再解释它们为何经常合在一起学习。

概率论的典型问题是：如果模型及其条件成立，某种结果发生的概率是多少？
统计学的典型问题是：已经取得这些数据，可以对未知对象作出什么判断，这一判断有多可靠？
因此，可以粗略地把前者理解为“从模型推导数据可能的表现”，
把\keyterm{统计推断}（statistical inference）理解为“利用数据认识模型及其未知部分”。
这只是帮助入门的概括，不是两个学科的严格边界，也不意味着统计推断是概率计算的简单逆运算。

\section{同一个抛硬币问题的两种研究方向}
\label{sec:prob-stat-coin}

考虑重复抛掷同一枚硬币。假设各次抛掷相互独立，且每次正面朝上的概率都为 $p$，
其中 $0\leq p\leq1$。这些是研究所采用的\keyterm{模型假设}，不能仅凭“抛的是同一枚硬币”就认为它们必然成立。
对第 $i$ 次抛掷，令随机变量 $X_i$ 在正面朝上时等于 $1$，反面朝上时等于 $0$，于是
\[
  \mathbb P(X_i=1)=p,\qquad \mathbb P(X_i=0)=1-p.
\]
这里 $\mathbb P$ 表示概率，大写 $X_i$ 表示随机变量，小写 $x_i$ 表示实际观测到的取值。

\begin{example}[给定概率，推导结果的可能性]
\label{ex:prob-stat-forward}
记 $S=X_1+\cdots+X_{100}$ 为抛掷 $100$ 次时正面出现的总次数。
如果 $p$ 已知，就可以计算“恰好出现 $60$ 次正面”的概率：
\begin{equation}
  \mathbb P(S=60)=\binom{100}{60}p^{60}(1-p)^{40}.
  \label{eq:prob-stat-binomial}
\end{equation}
其中 $\binom{100}{60}=100!/(60!\,40!)$ 表示从 $100$ 个位置中选出 $60$ 个正面位置的方式数，
符号 $n!$ 表示从 $1$ 到正整数 $n$ 的连乘积。
对每一种固定排列，由独立性可知其概率是 $p^{60}(1-p)^{40}$；
不同排列对应互不相容的事件，将它们的概率相加就得到式 \eqref{eq:prob-stat-binomial}。
若硬币正反面等可能，即 $p=1/2$，则该概率为 $\binom{100}{60}2^{-100}$。
这里的任务是根据模型计算事件概率，属于典型的概率论问题。
\end{example}

\begin{example}[给定观测，估计未知的概率]
\label{ex:prob-stat-inference}
现在假设 $p$ 未知，实际抛掷 $100$ 次，观察到 $60$ 次正面。
一种自然的\keyterm{参数估计}（parameter estimation）方法，是用样本中正面出现的频率估计 $p$。
在观测之前，所采用的\keyterm{估计量}（estimator）是随机变量
\begin{equation}
  \widehat p=\frac{1}{100}\sum_{i=1}^{100}X_i=\frac{S}{100};
  \label{eq:prob-stat-estimator}
\end{equation}
在取得这组数据后，对应的\keyterm{估计值}（estimate）为
\[
  \widehat p_{\mathrm{obs}}=\frac{1}{100}\sum_{i=1}^{100}x_i=\frac{60}{100}=0.6.
\]
估计量描述“遇到不同样本时怎样计算”，估计值则是这一规则在当前样本上得到的具体数值。
在这里采用的模型中，$p$ 是固定但未知的参数，样本及估计量的随机性来自重复抛掷。
\end{example}

观察到正面频率为 $0.6$，不能直接断言真实的 $p$ 就等于 $0.6$。
即使 $p=1/2$，例 \ref{ex:prob-stat-forward} 中的概率仍然大于零，
所以“正面次数不恰好等于 $50$”本身不能证明硬币不公平。
进一步的统计问题包括：估计可能有多大误差，以及数据是否对假设 $p=1/2$ 提供了足够的反对证据。
后一个问题属于\keyterm{假设检验}（hypothesis testing）；
它需要规定检验规则并考虑随机波动，不能仅凭估计值与假设值不完全相等就作出结论。
本例只说明问题的性质，不据此给出具体检验结论。

\section{二者的联系及常见误解}
\label{sec:prob-stat-connection}

概率论为统计推断提供数学基础。
例如，要判断式 \eqref{eq:prob-stat-estimator} 中的估计量是否可靠，就需要研究它在重复抽样中的分布和波动；
这些都是概率论能够处理的问题。
反过来，统计学通过数据估计未知参数、比较候选模型，并考察模型与研究问题是否相适应。
因此，在实际研究中，“建立模型、推导数据表现、收集数据、进行推断和检查模型”常常相互配合，
而不是只沿一个方向进行一次计算。

统计学也不只包含统计推断。
\keyterm{描述统计}（descriptive statistics）通过频数、平均数、图表等概括已观测的数据；
抽样和实验设计则关心应当如何取得有助于回答研究问题的数据。
例如，“这 $100$ 次抛掷中正面占 $60\%$”是在描述当前样本，
“据此估计每次抛掷正面朝上的概率约为 $0.6$”则已经涉及从样本推断未知参数。
后者的可靠性还依赖采集过程和模型条件，不能由前者自动保证。

课程名称中的\keyterm{数理统计}（mathematical statistics）着重研究统计方法的数学原理及其性质，
是统计学的重要组成部分，不能简单代替统计学的全部内容。
“概率论与数理统计”将两者合并讲授，主要体现知识依赖：先掌握随机变量、分布等概率基础，
再用这些工具研究参数估计和假设检验。
学习时应始终分清三个层次：模型作了什么假设，数据实际记录了什么，以及在这些条件下能够得到什么结论。

\end{document}
